\section{Occupational Characteristics and Heterogeneity}\label{app:characteristics}

\subsection{Inputs, timing, and common-support comparisons}

The characteristic exercise uses one variable at a time before considering a parsimonious joint model.  Every row first re-estimates the unaugmented exposure model on exactly the support available for that characteristic.  It then adds the characteristic interacted with young status and the post indicator, reuses the same occupation multipliers, and forms the augmented-minus-baseline difference within each draw.  Sample loss and coefficient movement are therefore distinct objects.  The corresponding paired minimum detectable effect is the two-sided, five-percent normal-theory quantity $(1.96+0.84)SE(\Delta)$; it describes resolution and is never used as a second rejection rule.

The leading computer-use variable is O*NET item 4.A.3.b.1, ``Interacting With Computers,'' from release 24.3 (May 2020).  Remotability is the Dingel--Neiman occupation classification.  Wage and observed education composition are calculated from the CPS preperiod, while externally coded education requirements are kept conceptually separate.  Routine and manual indices use preperiod O*NET inputs.  All scales, source files, checksums, weighted correlations, and support counts are stored in the characteristic result directory.

The primary matched comparison uses all 455 occupations with computer-use
ratings and crosses computer conditioning with family-month paths.  The four
Q5-by-post coefficients are $-0.09575$ without either adjustment, $-0.19660$
with computer use only, $0.03539$ with family paths only, and $-0.04665$ with
both.  Relative to the pooled baseline, the corresponding paired movements
are $-0.10086$ ($[-0.17723,-0.02448]$), $0.13114$
($[0.03517,0.22710]$), and $0.04910$ ($[-0.06418,0.16237]$).
The combined-minus-computer movement is $0.14995$
($[0.06039,0.23951]$), while combined-minus-family is $-0.08204$
($[-0.15461,-0.00947]$) under occupation clustering.  These distinct
conditional projections are not additive components.

The computer-use-by-post coefficient itself is $0.04602$ per weighted standard
deviation in the pooled computer-only model, with occupation interval
$[0.01171,0.08032]$, and $0.04868$ ($[0.01053,0.08684]$) in the combined
family-month model.  The weighted standard deviation is 0.97290 O*NET points,
and the beta--computer correlation is 0.7970.  Adding computer use reduces
Q5 information retention from 69.1 to 29.0 percent in the pooled model and
from 42.5 to 30.0 percent in the family-month model.  Thus coefficient
movement and inflation occur in a thin residual comparison; the computer
coefficient does not make computer use exogenous or purify the exposure
coefficient as AI.

I also evaluated formal omitted-confounder sensitivity rather than inserting
the grouped-binomial pseudo-$R^2$ into a linear formula.  The coefficient-
stability calculations of \citet{oster2019selection}, the partial-$R^2$
framework of \citet{cinelli2020sensitivity}, and the endogenous-control
selection model of \citet{diegert2022assessing} are derived for linear
regression targets with specific strength restrictions.  The present target is
one component of a vector-valued grouped-binomial fixed-effect projection with
generated exposure groups; no verified mapping supplies those scalar inputs.
A linear-probability companion would answer a different estimand.  I therefore
report no arbitrary breakdown value.  This applicability ruling is not
evidence that unobserved confounding is weak.

The separate common-characteristic block uses 347 occupations and holds its
support and primary-universe assignments fixed.  Relative to its $-0.12372$
baseline, computer use moves the coefficient by $-0.12475$
($[-0.20934,-0.04016]$).  The paired movements for remotability, wages,
education requirements, routine intensity, and manual/physical content each
include zero.  Adding all six gives $-0.23331$ and a paired movement of
$-0.10958$ ($[-0.22069,0.00152]$).  Family paths alone give $-0.00348$;
adding all six characteristics to that model gives $-0.10506$, with a paired
movement of $-0.10158$ ($[-0.20057,-0.00258]$).  The common-support block
covers 76.78 percent of rebuilt treatment-construction stock and likewise
does not identify a purified AI contrast.

\subsection{Expanded support-specific and cumulative audit}

The separate characteristic module in Tables~\ref{tab:appendix_characteristic_estimates} and \ref{tab:appendix_characteristic_information} reports the support-specific comparisons and the uncontaminated portion of the cumulative sequence.  Its one-at-a-time rows use each characteristic's maximal finite support, whereas its corrected joint rows use the literal 347-occupation common support.  An audit found that an earlier version of this module had retained historical full-static-panel exposure groups and a 341-occupation historical intersection.  The reported module now authenticates and uses the canonical corrected-preperiod membership and Webb normalization throughout; its 468-occupation baseline exactly reproduces $-0.1321$.  The support-specific computer comparison retains 455 occupations and holds those rebuilt primary-universe groups fixed, while the focused exercise requires both computer-use and remotability data and rebuilds groups on its 408-occupation support.  Their coefficients therefore differ by design, not by transcription.  Rows at and after the superseded unbounded shortfall control are removed rather than mixed with the corrected shortfall population.  Every addition and this corrective rerun are labeled post-outcome exploratory.

\begin{table}[!htbp]
\centering
\caption{Post-outcome exploratory characteristic-conditioning estimates}\label{tab:appendix_characteristic_estimates}
\scriptsize
\setlength{\tabcolsep}{3pt}
\resizebox{\textwidth}{!}{%
\begin{tabular}{lrrrrr}
\toprule
Conditioning block & Occupations & Baseline coefficient & Augmented coefficient [95\% CI] & Paired change [95\% CI] & Paired $\mathrm{MDE}_{80}$ \\
\midrule
\multicolumn{6}{l}{\textit{Panel A. One-at-a-time maximal-support comparisons}} \\
Computer use & 455 & $-0.096$ & $-0.197\;[-0.316,-0.077]$ & $-0.101\;[-0.177,-0.024]$ & 0.108 \\
Remotability & 408 & $-0.115$ & $-0.122\;[-0.222,-0.022]$ & $-0.007\;[-0.048,0.033]$ & 0.058 \\
Preperiod wage & 457 & $-0.144$ & $-0.161\;[-0.248,-0.075]$ & $-0.017\;[-0.044,0.009]$ & 0.038 \\
Education requirement & 404 & $-0.100$ & $-0.105\;[-0.203,-0.008]$ & $-0.005\;[-0.027,0.016]$ & 0.031 \\
Routine-task intensity & 449 & $-0.137$ & $-0.129\;[-0.221,-0.036]$ & $0.008\;[-0.011,0.027]$ & 0.027 \\
Manual/physical content & 455 & $-0.140$ & $-0.142\;[-0.289,0.004]$ & $-0.002\;[-0.092,0.089]$ & 0.130 \\
SOC2 by young by post & 468 & $-0.132$ & $-0.022\;[-0.162,0.119]$ & $0.111\;[0.008,0.213]$ & 0.146 \\
\addlinespace
\multicolumn{6}{l}{\textit{Panel B. One-at-a-time and joint comparisons on the 347-occupation common support}} \\
Computer use & 347 & $-0.124$ & $-0.248\;[-0.391,-0.106]$ & $-0.125\;[-0.209,-0.040]$ & 0.124 \\
Remotability & 347 & $-0.124$ & $-0.137\;[-0.239,-0.034]$ & $-0.013\;[-0.052,0.026]$ & 0.057 \\
Preperiod wage & 347 & $-0.124$ & $-0.135\;[-0.235,-0.036]$ & $-0.012\;[-0.029,0.006]$ & 0.025 \\
Education requirement & 347 & $-0.124$ & $-0.134\;[-0.235,-0.032]$ & $-0.010\;[-0.027,0.006]$ & 0.024 \\
Routine-task intensity & 347 & $-0.124$ & $-0.112\;[-0.220,-0.003]$ & $0.012\;[-0.024,0.048]$ & 0.051 \\
Manual/physical content & 347 & $-0.124$ & $-0.122\;[-0.255,0.010]$ & $0.001\;[-0.079,0.082]$ & 0.117 \\
All six characteristics & 347 & $-0.124$ & $-0.233\;[-0.388,-0.078]$ & $-0.110\;[-0.221,0.002]$ & 0.164 \\
\bottomrule
\end{tabular}}
\tabnote{All entries are post-outcome exploratory and were designed after the initial outcome pattern was known.  Each baseline--augmented pair uses identical occupations, the canonical rebuilt corrected-preperiod beta groups and Webb normalization, and 9,999 common occupation-level Rademacher multipliers.  Panel A uses each characteristic's maximal finite subset; Panel B uses the current 347-occupation common support, which covers 76.78 percent of primary-support January 2017--November 2022 treatment-construction stock.  For comparison, the unaugmented coefficient on all 468 occupations is $-0.132$, while the unaugmented coefficient after imposing the 347-occupation common support is $-0.124$; that support-only movement is kept separate from every paired conditioning change.  Rows from the superseded unbounded pandemic-shortfall construction are omitted; the corrected shortfall analysis uses a separate 363-occupation population and appears in Table~\ref{tab:appendix_shortfall}.  A confidence interval containing zero is a nondetection, not evidence of equivalence; conditioning does not identify a causal AI component.}
\end{table}


\begin{table}[!htbp]
\centering
\caption{Post-outcome exploratory support and target-information diagnostics}\label{tab:appendix_characteristic_information}
\scriptsize
\setlength{\tabcolsep}{3pt}
\resizebox{\textwidth}{!}{%
\begin{tabular}{lrrrrr}
\toprule
Conditioning block & Employment coverage (\%) & Information retained (\%) & Target VIF-like ratio & Effective occupation count & Top-five information share (\%) \\
\midrule
\multicolumn{6}{l}{\textit{Panel A. One-at-a-time maximal-support comparisons}} \\
Computer use & 96.75 & 29.02 & 3.45 & 52.58 & 22.34 \\
Remotability & 88.58 & 52.66 & 1.90 & 29.25 & 30.65 \\
Preperiod wage & 97.93 & 59.45 & 1.68 & 37.86 & 26.32 \\
Education requirement & 84.54 & 63.60 & 1.57 & 27.48 & 31.24 \\
Routine-task intensity & 98.45 & 65.19 & 1.53 & 43.26 & 24.59 \\
Manual/physical content & 96.85 & 30.61 & 3.27 & 51.05 & 22.34 \\
SOC2 by young by post & 100.00 & 20.81 & 4.81 & 53.06 & 21.87 \\
\addlinespace
\multicolumn{6}{l}{\textit{Panel B. Cumulative comparisons on the 341-occupation common support}} \\
Computer use & 76.77 & 29.33 & 3.41 & 48.88 & 21.98 \\
$+$ remotability & 76.77 & 28.46 & 3.51 & 44.97 & 24.72 \\
$+$ wage, education, routine, and manual & 76.77 & 21.08 & 4.74 & 58.00 & 20.28 \\
\bottomrule
\end{tabular}}
\tabnote{All diagnostics are post-outcome exploratory.  Employment coverage is the share of primary-support January 2017--November 2022 corrected treatment-construction stock retained.  ``Information retained'' is nuisance-adjusted target information divided by fixed-effect-adjusted raw target information; its reciprocal is the reported target VIF-like ratio.  These are design-information diagnostics, not an estimate of latent exposure reliability.  The effective occupation count is the inverse concentration of occupation-level conditional target-information shares, and the final column is the share supplied by the five largest occupations.  In the 341-occupation common-support sample, the 2017--2019-employment-weighted correlations of beta exposure with computer use, remotability, wage, education requirement, routine-task intensity, and manual/physical content are 0.807, 0.562, 0.429, 0.421, 0.197, and $-0.748$, respectively.  Rows based on the superseded unbounded pandemic-shortfall construction are omitted.  Substantial information loss makes residual coefficients less precise and does not purify an AI treatment.}
\end{table}


\subsection{Pandemic shortfall}

The corrected exercise uses ages 22--65 throughout and admits only one-to-one occupation mappings with sufficient 2017--2019 histories.  It retains 363 occupations, rather than the central 468.  Linear trend predictions for total employment are bounded below at zero; predicted young employment shares are bounded to $[0,1]$.  This is consequential: before bounding there are 174 negative total predictions and 1,258 out-of-range share predictions.  The characteristics are the average bounded trend-prediction-minus-realized gaps from January 2020 through November 2022; December 2022 remains the transition month.

On this narrower common support, the pooled and family-month baseline coefficients are $-0.16417$ and $-0.08708$.  Total-stock conditioning moves them by 0.00040 and $-0.00050$; young-relative conditioning moves them by 0.00114 and 0.00613.  Table~\ref{tab:appendix_shortfall} reports linked-household sampling sensitivities that regenerate the shortfall inside each draw.  All four conditioning-movement intervals include zero.  Jointly regenerating exposure labels changes an average of 1.99 occupations per draw and leaves the young-relative family-month movement at 0.00613 with sampling-sensitivity standard error 0.00695.

All 399 linked-household draws finish with zero failed joint refits.  The
largest binding bootstrap-of-bootstrap endpoint Monte Carlo standard error is
0.00979 log point, below the declared 0.01 threshold, so the predeclared
expansion to 799 draws is not triggered.

The controls are generated from related CPS outcomes and may share sampling error and regression-to-the-mean with the dependent stocks.  A deterministic household split gives a similar average but heterogeneous directions; it does not eliminate PSU-level common error.  The bounded functional form and changed population also remain consequential.  The exercise therefore does not establish that pandemic dynamics are absent, and the linked-household intervals are not complete CPS design-based inference.

\begin{table}[!htbp]
\centering
\caption{Corrected bounded pandemic-shortfall diagnostics}\label{tab:appendix_shortfall}
\resizebox{\textwidth}{!}{%
\begin{tabular}{llrrr}
\toprule
Shortfall & Model & Baseline Q5--Q1 & Conditioned Q5--Q1 & Movement [sampling interval] \\
\midrule
Total stock & Pooled & $-0.16417$ & $-0.16377$ & $0.00040\;[-0.00655,0.00661]$ \\
Total stock & SOC2 $\times$ month & $-0.08708$ & $-0.08758$ & $-0.00050\;[-0.01377,0.00970]$ \\
Young relative & Pooled & $-0.16417$ & $-0.16303$ & $0.00114\;[-0.00509,0.00573]$ \\
Young relative & SOC2 $\times$ month & $-0.08708$ & $-0.08095$ & $0.00613\;[-0.00889,0.01835]$ \\
\bottomrule
\end{tabular}}
\tabnote{The common population contains 363 one-to-one-mapped occupations with sufficient preperiod histories.  Total-employment predictions are bounded below at zero and young-share predictions to $[0,1]$.  Movement is the conditioned minus baseline exposure coefficient.  Brackets are basic intervals from 399 linked-household draws that regenerate the shortfall but hold exposure labels fixed.  They are released-weight sampling sensitivities, not occupation-wild or complete CPS design-based intervals, and are not mechanically combined with the article's shock-cluster variance.  Joint label regeneration gives the same displayed point estimates and very similar uncertainty.}
\end{table}


\subsection{Industry, education, and exact age}

The industry analysis is rebuilt at the occupation-by-broad-industry-by-age-by-month level rather than assigning each occupation a dominant industry.  The risk set covers all 468 occupations and 3,127 preperiod-connected occupation--industry cells.  Relative to the stratified baseline coefficient of $-0.1376$, adding 13 broad-industry-specific young-by-post shifts gives $-0.0987$.  The paired movement is $0.0389$, with interval $[0.0082,0.0697]$ and paired $\mathrm{MDE}_{80}=0.0444$.  The conditioned coefficient remains below zero.  These are statements about the microcell objective, not firm-level adoption.

Education-specific models are estimated on a common occupational support.  The BA-or-higher estimate is $-0.1380$ and the non-BA estimate is $-0.1226$; their paired difference is $-0.0154$, with interval $[-0.1557,0.1250]$ and paired $\mathrm{MDE}_{80}=0.2014$.  The comparison is too imprecise to distinguish moderate education differences.  Exact-age estimates for ages 22, 23, 24, and 25 are reported as one jointly estimated exploratory collection with simultaneous intervals.  No single age is selected after inspection.

\begin{table}[!htbp]
\centering
\caption{Industry, education, and exact-age comparisons}\label{tab:app_heterogeneity}
\resizebox{\textwidth}{!}{%
\begin{tabular}{lrrrr}
\toprule
Model & Occupations & Coefficient & Pointwise 95\% interval & Simultaneous 95\% interval \\
\midrule
Occupation--industry-cell baseline & 468 & $-0.138$ & $[-0.228,-0.048]$ & --- \\
Industry-conditioned & 468 & $-0.099$ & $[-0.188,-0.009]$ & --- \\
BA or higher & 429 & $-0.138$ & $[-0.275,-0.001]$ & $[-0.293,0.017]$ \\
No BA & 429 & $-0.123$ & $[-0.231,-0.014]$ & $[-0.243,-0.002]$ \\
Age 22 & 439 & $-0.093$ & $[-0.212,0.026]$ & $[-0.239,0.053]$ \\
Age 23 & 439 & $-0.140$ & $[-0.252,-0.028]$ & $[-0.280,-0.0003]$ \\
Age 24 & 439 & $-0.222$ & $[-0.358,-0.086]$ & $[-0.394,-0.050]$ \\
Age 25 & 439 & $-0.082$ & $[-0.188,0.024]$ & $[-0.212,0.047]$ \\
\bottomrule
\end{tabular}}
\tabnote{Both industry rows use occupation--broad-industry--age--month weighted-stock cells.  The baseline includes common Q2--Q5 and Webb interactions; the conditioned row additionally includes 13 broad-industry-specific young-by-post shifts.  Education and age collections use their own common occupational support and 9,999 common occupation multipliers.  The four-age common-score equality test gives $\chi^2(3)=8.21$, $p=.0419$, but every simultaneous age-minus-pooled interval includes zero; no individual age departure is promoted.}
\end{table}


\subsection{Older comparison bands, age composition, and enrollment}

Table~\ref{tab:app_cohort_enrollment} varies the older comparison group while
holding the current exposure labels and occupational support fixed.  The pooled
coefficient remains negative across all five bands, but its magnitude is not
constant.  Against the 26--65 benchmark, paired movements for the 26--30,
31--40, and 41--50 comparisons have occupation-multiplier intervals
$[-0.0412,0.0644]$, $[-0.0610,0.0056]$, and $[-0.0413,0.0271]$.
The 51--65 movement is 0.0350 ($[0.0036,0.0663]$).  Under family-month
conditioning, every coefficient interval and all four corresponding paired
movement intervals include zero when shocks are assigned at the occupation
level.  Under 22-family shock assignment the ages 22--25 versus 26--30 paired
movement instead has interval $[-0.1712,-0.0041]$; the two procedures answer
different stochastic questions and are not combined, and neither is evidence
of a causal age effect.

Fixing the older group's exact-age composition at its preperiod shares gives a
pooled coefficient of $-0.1313$ ($[-0.2202,-0.0425]$) and a family-month
coefficient of $-0.0218$ ($[-0.1597,0.1161]$).  Relative to observed age
composition, the paired movements are 0.0008 ($[-0.0011,0.0027]$) and
$-0.0001$ ($[-0.0026,0.0023]$).

Enrollment restrictions change the numerator rather than assigning exposure to
nonemployed people.  With all ages 26--65 in the denominator, restricting ages
22--25 to nonenrolled workers changes the pooled coefficient from $-0.1321$ to
$-0.1239$; the paired movement is 0.0082 ($[-0.0125,0.0289]$).  On the common
22--54 enrollment-question universe, restricting both age groups to nonenrolled
workers changes the pooled coefficient from $-0.1407$ to $-0.1354$, a paired
movement of 0.0053 ($[-0.0144,0.0250]$).  The analogous family-month movements
are $-0.0068$ ($[-0.0393,0.0257]$) and $-0.0059$
($[-0.0377,0.0260]$).  Enrollment may itself respond to labor-market conditions,
so these comparisons describe composition rather than a causal correction.
Older workers throughout are comparison groups, not untreated controls.

\begin{table}[!htbp]
\centering
\caption{Older comparison bands and enrollment restrictions}\label{tab:app_cohort_enrollment}
\resizebox{\textwidth}{!}{%
\begin{tabular}{lrr}
\toprule
Comparison & Pooled coefficient and 95\% interval & Family-month coefficient and 95\% interval \\
\midrule
Ages 22--25 versus 26--65 & $-0.1321\;[-0.2211,-0.0431]$ & $-0.0217\;[-0.1607,0.1174]$ \\
Ages 22--25 versus 26--30 & $-0.1205\;[-0.1881,-0.0530]$ & $-0.1093\;[-0.2328,0.0141]$ \\
Ages 22--25 versus 31--40 & $-0.1598\;[-0.2455,-0.0742]$ & $-0.0247\;[-0.1709,0.1216]$ \\
Ages 22--25 versus 41--50 & $-0.1392\;[-0.2344,-0.0440]$ & $-0.0022\;[-0.1429,0.1385]$ \\
Ages 22--25 versus 51--65 & $-0.0971\;[-0.2037,0.0094]$ & $-0.0021\;[-0.1627,0.1584]$ \\
Older exact-age composition fixed & $-0.1313\;[-0.2202,-0.0425]$ & $-0.0218\;[-0.1597,0.1161]$ \\
Young numerator nonenrolled & $-0.1239\;[-0.2154,-0.0324]$ & $-0.0285\;[-0.1666,0.1097]$ \\
Both groups nonenrolled, ages 22--54 & $-0.1354\;[-0.2243,-0.0465]$ & $-0.0356\;[-0.1704,0.0991]$ \\
\bottomrule
\end{tabular}}
\tabnote{Intervals use 9,999 occupation multipliers.  The fixed-composition row holds the 26--65 exact-age shares at their preperiod values.  The young-numerator row retains all employed ages 26--65 in the denominator.  The last row restricts both groups to $\mathrm{SCHLCOLL}=5$ on the common ages 22--54 question universe.  Older workers are comparison groups, not untreated controls; an interval containing zero indicates non-detection rather than equivalence.}
\end{table}


\subsection{Young-worker composition and national cohort supply}

Table~\ref{tab:app_young_composition} keeps employment-conditioned composition
separate from national cohort supply.  The Q1--Q5 columns summarize employed
ages 22--25 in occupations on the fixed 468-occupation support.  The national
profile uses every positive-weight CPS record ages 22--25 and assigns no
occupation or exposure quintile.  Across Q1--Q5, the pre-to-post enrolled-share
changes are $-0.10$, $-0.74$, $-2.00$, $-1.46$, and $-1.80$ percentage points;
BA-or-higher shares change by 1.79, 1.71, 2.69, 1.48, and 3.81 points.  No
employed-quintile mean-age change exceeds 0.017 year in absolute value, and the
largest exact-age-share change is 1.62 points.  Nationally, enrollment changes
from 0.2336 to 0.2211, BA-or-higher from 0.2941 to 0.3121, employment from
0.7110 to 0.7247, and mean age by $-0.0085$ year; average monthly weighted
population rises from 16.94 to 17.60 million.  These are descriptive
composition and selection patterns.  Current enrollment may respond to
labor-market conditions, and the profiles do not identify a cohort correction
or assign exposure to nonworkers.

\begin{table}[!htbp]
\centering
\caption{Education, enrollment, and age composition of ages 22--25}\label{tab:app_young_composition}
\textit{Panel A: Annual enrollment and BA-or-higher shares}\par\smallskip
\resizebox{\textwidth}{!}{%
\begin{tabular}{lrrrrrrr}
\toprule
Year & Q1 & Q2 & Q3 & Q4 & Q5 & National population & National population (millions/month) \\
\midrule
2017 & .135/.085 & .201/.239 & .191/.397 & .142/.473 & .171/.477 & .244/.275 & 17.376 \\
2018 & .113/.089 & .207/.232 & .188/.439 & .149/.470 & .171/.484 & .239/.281 & 17.295 \\
2019 & .121/.080 & .208/.238 & .168/.449 & .151/.489 & .162/.518 & .236/.293 & 16.865 \\
2020 & .121/.098 & .190/.247 & .173/.463 & .125/.517 & .158/.562 & .232/.314 & 16.848 \\
2021 & .118/.101 & .177/.245 & .167/.444 & .149/.474 & .164/.555 & .223/.301 & 16.541 \\
2022 & .120/.088 & .186/.253 & .159/.459 & .130/.489 & .152/.538 & .225/.302 & 16.673 \\
2023 & .118/.110 & .181/.245 & .154/.470 & .123/.480 & .152/.534 & .219/.306 & 17.659 \\
2024 & .121/.109 & .195/.249 & .147/.469 & .129/.517 & .136/.559 & .221/.315 & 17.369 \\
2025 & .121/.113 & .189/.269 & .157/.463 & .128/.493 & .151/.580 & .222/.316 & 17.638 \\
2026 & .124/.096 & .186/.283 & .164/.470 & .127/.516 & .140/.575 & .222/.313 & 17.815 \\
\bottomrule
\end{tabular}}

\medskip
\textit{Panel B: Preperiod/postperiod age composition}\par\smallskip
\resizebox{\textwidth}{!}{%
\begin{tabular}{lrrrrrr}
\toprule
Group & Mean age & YEAR--age proxy & Age 22 & Age 23 & Age 24 & Age 25 \\
\midrule
Q1 & 23.482/23.486 & 1995.9/2000.8 & .257/.251 & .251/.250 & .245/.261 & .247/.238 \\
Q2 & 23.481/23.469 & 1995.9/2000.9 & .256/.253 & .253/.257 & .244/.258 & .246/.232 \\
Q3 & 23.642/23.638 & 1995.8/2000.7 & .198/.200 & .248/.241 & .267/.279 & .286/.280 \\
Q4 & 23.690/23.681 & 1995.8/2000.6 & .190/.184 & .235/.245 & .270/.276 & .305/.295 \\
Q5 & 23.657/23.674 & 1995.8/2000.6 & .194/.192 & .241/.236 & .280/.279 & .286/.293 \\
National & 23.514/23.506 & 1995.9/2000.8 & .247/.246 & .248/.249 & .250/.258 & .255/.247 \\
\bottomrule
\end{tabular}}
\tabnote{Q1--Q5 condition on employment and the fixed 468-occupation support.  ``National population'' uses every positive-\texttt{WTFINL} CPS record ages 22--25 and assigns neither an occupation nor an exposure quintile to nonworkers.  Entries in Panel A before the final column are enrollment share/BA-or-higher share; entries in Panel B are preperiod/postperiod.  The preperiod is January 2017--November 2022 and the postperiod January 2023--July 2026; December 2022 is omitted, October 2025 is absent, and 2026 is partial.  Annual national population is the person-weight sum divided by observed months.  YEAR minus age is an approximate birth-year proxy, not an age--period--cohort coefficient.  Complete annual age/cohort cells remain in the machine-readable composition files.}
\end{table}

